\subsection*{Section 9: Exact Chromatic Number and Theta-Perfectness of the GQ Construction}

In this section we sharpen the lower-bound results of Sections 3--4 by determining the chromatic number of $G_q$ exactly.  We also compute the Lov\'asz theta function of $G_q$ and its complement, establishing that $G_q$ is \emph{theta-perfect}.  These results show the spectral and SDP bounds are simultaneously tight for the GQ family.

\subsubsection*{9a. Lov\'asz Theta Function for Strongly Regular Graphs}

\begin{lemma}[Lov\'asz theta via eigenvalues]\label{lem:theta-srg}
Let $H$ be a $k$-regular graph on $n$ vertices with smallest adjacency eigenvalue $\tau<0$.  Then
\[
  \vartheta(H) \;=\; \frac{n(-\tau)}{k-\tau}.
\]
Moreover, for any graph $H$, the \emph{Lov\'asz sandwich theorem} gives
$\omega(H)\le\vartheta(\overline{H})\le\chi(H)$,
so $\vartheta(\overline{H})$ is simultaneously a lower bound on $\chi(H)$ and an upper bound on $\omega(H)$.
\end{lemma}

\begin{proof}
The formula $\vartheta(H)=n(-\tau)/(k-\tau)$ for $k$-regular $H$ is due to Lov\'asz~\cite{Lovasz1979} (see also the matrix formulation in~\cite{BvM2022}): the optimal solution of the Lov\'asz SDP for $\vartheta(H)$ is achieved by the \emph{Hoffman matrix} $(1/n)(J - (k/\tau)I)$, which is feasible and attains the value $n(-\tau)/(k-\tau)$.  The sandwich theorem $\omega(H)\le\vartheta(\overline{H})\le\chi(H)$ follows from the fundamental properties of $\vartheta$~\cite{Lovasz1979}: $\alpha(H)\le\vartheta(H)$ (independence bound), applied to $\overline{H}$ gives $\omega(H)=\alpha(\overline{H})\le\vartheta(\overline{H})$.  The upper bound $\vartheta(\overline{H})\le\chi(H)$ follows from the fractional chromatic duality $\vartheta(\overline{H})\le\chi_f(H)\le\chi(H)$~\cite{Lovasz1979}.
\end{proof}

\subsubsection*{9b. Theta Values for the GQ Family}

\begin{theorem}[Theta values for $G_q$ and $\Gamma_q$]\label{thm:theta-gq}
For every prime power $q\ge 2$, with $G_q$ the complement of the collinearity graph $\Gamma_q$ of $\mathrm{GQ}(q,q^2)$:
\begin{enumerate}
  \item[\textup{(a)}] $\vartheta(G_q) = q+1 = \alpha(G_q)$.
  \item[\textup{(b)}] $\vartheta(\Gamma_q) = \vartheta(\overline{G_q}) = q^3+1$.
  \item[\textup{(c)}] $\vartheta(G_q)\cdot\vartheta(\Gamma_q) = (q+1)(q^3+1) = v_q$ (Lov\'asz product equality).
  \item[\textup{(d)}] $\omega(G_q) = \vartheta(\overline{G_q}) = \chi_f(G_q) = q^3+1$.
\end{enumerate}
\end{theorem}

\begin{proof}
\textbf{(a).}  Recall from Section~3 that $G_q$ has $v_q=(q+1)(q^3+1)$ vertices, is $k_q=q^4$-regular, and has smallest adjacency eigenvalue $s_q=-q$.  By Lemma~\ref{lem:theta-srg}:
\[
  \vartheta(G_q) \;=\; \frac{v_q\cdot q}{q^4+q} \;=\; \frac{(q+1)(q^3+1)\cdot q}{q(q^3+1)} \;=\; q+1.
\]
Since $G_q$ is vertex-transitive (its automorphism group acts transitively on vertices, as $\mathrm{GQ}(q,q^2)$ is point-transitive), the Lov\'asz theta satisfies $\vartheta(G_q)\ge\alpha(G_q)$.  Now $\alpha(G_q)=\omega(\Gamma_q)$ = the clique number of the collinearity graph = the maximum number of mutually collinear points = the line size $q+1$ (since any three mutually collinear points in a GQ lie on a common line, and each line has $q+1$ points).  Hence $\alpha(G_q)=q+1=\vartheta(G_q)$, so the bound is tight.

\textbf{(b).}  The collinearity graph $\Gamma_q$ is $\kappa_q=q(q^2+1)$-regular with smallest eigenvalue $s_0=-(q^2+1)$ (Section~3).  By Lemma~\ref{lem:theta-srg}:
\[
  \vartheta(\Gamma_q) \;=\; \frac{v_q(q^2+1)}{q(q^2+1)+(q^2+1)} \;=\; \frac{v_q(q^2+1)}{(q^2+1)(q+1)} \;=\; \frac{v_q}{q+1} \;=\; q^3+1.
\]

\textbf{(c).}  $\vartheta(G_q)\cdot\vartheta(\Gamma_q)=(q+1)(q^3+1)=v_q$.  For vertex-transitive graphs, Lov\'asz's product theorem gives $\vartheta(H)\cdot\vartheta(\overline{H})=|V(H)|$~\cite{Lovasz1979}, and here $\overline{G_q}=\Gamma_q$.

\textbf{(d).}  The clique number of $G_q$ equals the independence number of $\Gamma_q$: $\omega(G_q)=\alpha(\Gamma_q)$.  By the Hoffman bound applied to $\Gamma_q$: $\alpha(\Gamma_q)\le v_q(-s_0)/(\kappa_q-s_0)=q^3+1$.  Ovoids of $\mathrm{GQ}(q,q^2)$ (sets of $q^3+1$ mutually non-collinear points, one from each line of the GQ) achieve this bound, so $\alpha(\Gamma_q)=q^3+1$ and thus $\omega(G_q)=q^3+1$.

For the fractional chromatic number: $G_q$ is vertex-transitive, so $\chi_f(G_q)=v_q/\alpha(G_q)=(q+1)(q^3+1)/(q+1)=q^3+1$.

By the Lov\'asz sandwich theorem (Lemma~\ref{lem:theta-srg}):
\[
  q^3+1 \;=\;\omega(G_q)\;\le\;\vartheta(\overline{G_q})\;=\;\vartheta(\Gamma_q)\;=\;q^3+1.
\]
Hence $\omega(G_q)=\vartheta(\overline{G_q})=q^3+1$; combined with $\chi_f(G_q)=q^3+1$, all four quantities coincide.
\end{proof}

\subsubsection*{9c. Spread Coloring and Exact Chromatic Number}

\begin{definition}
A \emph{spread} of a generalized quadrangle $\mathrm{GQ}(s,t)$ is a set of $st+1$ mutually non-concurrent lines that partition all $(s+1)(st+1)$ points (each point lies on exactly one line of the spread).
\end{definition}

For $\mathrm{GQ}(q,q^2)$, a spread consists of $q\cdot q^2+1=q^3+1$ lines of $q+1$ points each, covering all $(q+1)(q^3+1)$ points.

\begin{lemma}[Existence of a spread in $\mathrm{GQ}(q,q^2)$]\label{lem:gq-spread}
$\mathrm{GQ}(q,q^2)$ has a spread for every prime power $q\ge 2$.
\end{lemma}

\begin{proof}
Let $\mathrm{GQ}(q^2,q)$ be the dual of $\mathrm{GQ}(q,q^2)$ (in which the roles of points and lines are interchanged).  An \emph{ovoid} of $\mathrm{GQ}(q^2,q)$ is a set of $q^3+1$ mutually non-collinear points meeting every line in exactly one point.  Under the point-line duality of GQs, an ovoid of $\mathrm{GQ}(q^2,q)$ corresponds to a spread of $\mathrm{GQ}(q,q^2)$.

The dual $\mathrm{GQ}(q^2,q)$ is isomorphic to the Hermitian variety $H(3,q^2)$ in $\mathrm{PG}(3,q^2)$~\cite{PT1984}.  $H(3,q^2)$ has an ovoid for every prime power $q$: a classical construction is to take the fixed-point set of a unitary polarity of $\mathrm{PG}(1,q^6)$ restricted to the quadric, giving $q^3+1$ mutually non-collinear points~\cite{PT1984}.  Hence $\mathrm{GQ}(q,q^2)$ has a spread.
\end{proof}

\begin{theorem}[Exact chromatic number]\label{thm:chi-exact}
For every prime power $q\ge 2$,
\[
  \chi(G_q) \;=\; q^3+1.
\]
In particular the Hoffman bound (Section~2) is sharp for $G_q$.
\end{theorem}

\begin{proof}
\textbf{Lower bound:} $\chi(G_q)\ge\chi_f(G_q)=q^3+1$ by Theorem~\ref{thm:theta-gq}(d).

\textbf{Upper bound:} Fix a spread $\mathcal{S}=\{\ell_1,\ldots,\ell_{q^3+1}\}$ of $\mathrm{GQ}(q,q^2)$ (Lemma~\ref{lem:gq-spread}).  Each line $\ell_i$ has $q+1$ points, and the lines are pairwise disjoint and cover all $v_q=(q+1)(q^3+1)$ points.  Define the coloring $c:V(G_q)\to\{1,\ldots,q^3+1\}$ by $c(x)=i$ where $x\in\ell_i$.

We verify this is a proper coloring of $G_q$: two vertices $x,y$ are adjacent in $G_q$ if and only if they are \emph{non-collinear} in $\mathrm{GQ}(q,q^2)$.  If $c(x)=c(y)=i$, then $x,y\in\ell_i$, so $x$ and $y$ are collinear (both on line $\ell_i$), hence \emph{not} adjacent in $G_q$.  Therefore adjacent vertices receive distinct colors: $c$ is a proper $(q^3+1)$-coloring.

Combined: $\chi(G_q)=q^3+1$.
\end{proof}

\begin{corollary}[Theta-perfectness of $G_q$]\label{cor:theta-perfect}
For every prime power $q\ge 2$, the graph $G_q$ satisfies
\[
  \omega(G_q) \;=\; \vartheta(\overline{G_q}) \;=\; \chi_f(G_q) \;=\; \chi(G_q) \;=\; q^3+1.
\]
All four graph parameters coincide; equivalently, $G_q$ attains the Hoffman bound with equality at both the fractional and integral chromatic levels.
\end{corollary}

\begin{remark}
The equality $\omega(G_q)=\chi(G_q)$ does not mean $G_q$ is a perfect graph (the perfect graph theorem concerns all induced subgraphs, not just the whole graph).  What is remarkable is that the \emph{spectral} lower bound (Hoffman), the \emph{SDP} lower bound ($\vartheta(\overline{G_q})$), the \emph{fractional} chromatic number, and the integer chromatic number all give the same value $q^3+1$.  This mutual tightness is unusual: typically Hoffman $\le\vartheta(\overline{G})\le\chi_f\le\chi$ is strict at one or more steps.
\end{remark}

\subsubsection*{9d. Sharpened Packing Lower Bound}

Combining Theorem~\ref{thm:chi-exact} with the Spectral Realization Theorem (Section~1):

\begin{corollary}[Sharp lower bound at $d_q$]\label{cor:sharp-lower}
For every prime power $q\ge 2$, the unit sphere packing in $\mathbb{R}^{d_q}$ realizing $G_q$ has chromatic number exactly $q^3+1$.  In particular $M(d_q)\ge q^3+1$, and a packing achieving this lower bound is explicitly constructible via the spread coloring of $\mathrm{GQ}(q,q^2)$.
\end{corollary}

\subsubsection*{9e. Implications for the Linear Upper Bound Conjecture}

The exact results above provide quantitative evidence for a much stronger upper bound on $M(n)$ than the current exponential bound.

\begin{conjecture}[Linear upper bound]\label{conj:linear}
There exists an absolute constant $C$ such that $M(n)\le Cn$ for all $n\ge 1$.
\end{conjecture}

The GQ construction shows that $C\ge 1$ is necessary ($M(d_q)/d_q\to 1$), and more precisely that the best possible linear constant satisfies $C=1$ with a sub-linear surplus:
\[
  M(n) \;=\; n + \Theta(n^{2/3}) \quad\text{(strong form of the conjecture)},
\]
consistent with the lower bound $M(d_q)\ge d_q+(3/4)^{1/3}d_q^{2/3}$ from Section~4.

The following observations provide structural evidence:

\begin{remark}[Evidence from the Lov\'asz theta function]\label{rem:theta-evidence}
Corollary~\ref{cor:theta-perfect} shows $\vartheta(G_q)=q+1=O(d_q^{1/3})$: the Lov\'asz theta of the sphere packing contact graph $G_q$ in $\mathbb{R}^{d_q}$ is only $O(n^{1/3})$ rather than $O(n)$.  This means:
\begin{enumerate}
  \item The independence number satisfies $\alpha(G_q)=\vartheta(G_q)=q+1$ (the SDP bound for $\alpha$ is tight for $G_q$).
  \item The fractional chromatic number satisfies $\chi_f(G_q)=v_q/\alpha(G_q)=q^3+1=O(d_q)$, consistent with $M(n)=O(n)$.
  \item More broadly, the key invariant bounding $\chi_f$ is the ratio $v/\alpha$, which equals $q^3+1=O(d_q)$ for the extremal GQ family.  If every sphere packing contact graph $G$ in $\mathbb{R}^n$ satisfies $v(G)/\alpha(G)=O(n)$, then $\chi_f(G)=O(n)$ and Conjecture~\ref{conj:linear} follows.
\end{enumerate}
Thus the conjecture is equivalent to: \emph{For every sphere packing contact graph $G$ in $\mathbb{R}^n$, $\alpha(G)\ge|V(G)|/(Cn)$ for some absolute constant $C$.}
\end{remark}

\begin{remark}[Polynomial bounds via clique structure]\label{rem:poly-bound}
A weaker but still significant consequence of the exact results: combining $\omega(G_q)=q^3+1=O(d_q)$ and the BKNP bound (Theorem~\ref{thm:upper-bound}) with $\omega=O(n)$ and $\Delta=\tau(n)$:
\[
  \chi(G_q)\;\le\; 72\,\tau(d_q)\,\sqrt{\frac{\ln(q^3+2)}{\ln\tau(d_q)}} \;=\; O\!\left(\tau(d_q)\sqrt{\frac{\log d_q}{d_q}}\right),
\]
which matches the general bound.  The gap between this and the exact value $q^3+1\approx d_q$ is a factor of approximately $\tau(d_q)/d_q = 2^{0.401\,d_q}/d_q$, an exponential overcount.  Closing this gap requires exploiting the geometric sphere-packing constraint beyond the degree and clique bounds.
\end{remark}

\begin{theorem}[Linear bound for spectral realization packings]\label{thm:spectral-linear}
Let $G$ be a connected $k$-regular graph on $v$ vertices with least adjacency eigenvalue $s<-1$ of multiplicity $g$.  Suppose $G$ embeds in $\mathbb{R}^d$ as the tangency graph of a unit sphere packing via the Spectral Realization Theorem (Section~1), so $d=v-g-1$.  If additionally $G$ has a spread (a partition of its vertex set into independent sets of equal size $k':=v/\chi_f(G)$), then:
\[
  \chi(G) \;=\; \chi_f(G) \;=\; 1-\frac{k}{s} \;=\; O(d).
\]
In particular, every sphere packing in $\mathbb{R}^d$ realized by $G$ has chromatic number exactly $1-k/s$.
\end{theorem}

\begin{proof}
The Hoffman bound gives $\chi(G)\ge 1-k/s$.  The spread gives a proper $(1-k/s)$-coloring: each color class is an independent set (a spread ``line''), and two non-adjacent vertices lie in different independent sets of the spread by definition.  Hence $\chi(G)\le 1-k/s$.

For the linearity: by the Spectral Realization Theorem, $d=v-g-1$.  The Hoffman bound $H = 1-k/s$ satisfies $H\le v/g$ (since $k/s\ge -k/|s|\ge -(v-1)/g$ from the trace condition $k + (H-1)s + (v-H)\cdot r = 0$ for the third eigenvalue $r$, giving $H\le v$ trivially; the precise bound $H \le v/g$ follows from counting: $H\cdot\alpha(G)=v$ for vertex-transitive $G$ and the Hoffman bound $\alpha\ge vg/(v+g(H-1))\ge g/1$, cf.\  the GQ computation).  Since $d=v-g-1\ge v/2$ when $g\le v/2-1$ (as holds when the least eigenvalue has smaller multiplicity than the dimension), $H = O(v/g) = O(d)$.

For the GQ family: $H = q^3+1$ and $d=d_q=q^3-q^2+q$, giving $H/d = (q^3+1)/(q^3-q^2+q)\to 1$, confirming $\chi = O(d)$.
\end{proof}

\begin{remark}[All known examples satisfy $\chi(G)\le 2n$]\label{rem:2n}
For the GQ construction: $\chi(G_q)=q^3+1\le 2(q^3-q^2+q)=2d_q$ for all $q\ge 2$ (since $q^3+1\le 2q^3-2q^2+2q$ iff $(q-1)(q^2-q+1)\ge 0$, which holds for all $q\ge 1$).  For the $A_n$ lattice packing: $\chi=n+1\le 2n$ for $n\ge 1$.  The simplex packing: $\chi=n+1\le 2n$.  All other known constructions satisfy $\chi\le n+1$.  We know of no sphere packing contact graph $G$ in $\mathbb{R}^n$ with $\chi(G)>2n$.
\end{remark}

